\section{Discussion}

\subsection{Solvability Validation of Synthesized Natural Language Requirements}
\label{sec:human_eval}
To ensure the solvability of our synthesized natural language requirements and mitigate potential risks introduced by the automated pipeline, such as LLM hallucinations, ambiguity, or the omission of critical implementation details, we conducted a rigorous human solvability assurance process.

\paragraph{\textbf{Setup and Annotation Process.}}
We randomly sampled 30 tasks from the constructed benchmark, representing approximately 19\% of the total tasks. The evaluation employed a strict 3-point Likert scale to assess information completeness, as defined in \cref{tab:human_eval_criteria}. To ensure a professional and accurate assessment, we recruited two evaluators from the author team, both of whom are computer science researchers with over 2 years of Python development experience. Specifically, they examined whether the provided prompt contained all the necessary contextual information required to reproduce the functionality of the gold patch. To ensure the reliability of the evaluation results, evaluators were permitted to access the internet to clarify any unfamiliar technical concepts or domain-specific knowledge. All evaluators are paid according to the relevant policies\footnote{\url{https://www.worker.gov}} (i.e., \$7.5 per hour).

\begin{table}[t]
  \caption{Criteria for Human Evaluation on Information Completeness.}
  \label{tab:human_eval_criteria}
  \begin{tabular}{c p{0.8\linewidth}}
    \toprule
    \textbf{Score} & \textbf{Definition of Information Completeness} \\
    \midrule
    2 & \textbf{Fully Solvable:} The prompt perfectly covers all functional requirements and constraints mentioned in the PR, allowing for a correct implementation without ambiguity. \\
    \midrule
    1 & \textbf{Solvable with Inference:} The prompt covers the majority of feature aspects. It may omit minor, non-critical details, but the feature remains solvable through reasonable developer inference. \\
    \midrule
    0 & \textbf{Unsolvable:} The prompt is too narrow or broad, misses critical logical constraints, or looks for something different from the PR, rendering the task unsolvable. \\
    \bottomrule
  \end{tabular}
\end{table}

\paragraph{\textbf{Verification Results.}}
The manual review yielded highly positive results regarding FeatBench's quality. Of the 30 sampled instances, 28 (\textbf{93.3\%}) received a perfect score of 2, indicating that the synthesized requirements were comprehensive and unambiguous. The remaining 2 instances (\textbf{6.7\%}) received a score of 1, typically due to minor implicit context that a developer could reasonably infer. No instances were marked as Unsolvable (Score 0). The average score across the sample was \textbf{1.93} out of 2. These quantitative results confirm that our automated construction pipeline consistently generates high-quality, information-rich task requirements, establishing a solid foundation for subsequent model evaluation.

\subsection{Threats to Validity}
There are three main threats to the validity of our work:

\paragraph{\textbf{Quality of LLM-generated requirements.}}
A primary threat lies in whether the synthesized natural language requirements faithfully reflect the original developer's intent.
Since we utilize LLMs to reverse-engineer these requirements from code changes, there is a risk of hallucinations or ambiguity in the task descriptions.
To mitigate this, we employed a reconstructive prompting strategy: we explicitly instructed an LLM to craft descriptions sufficiently detailed to allow another LLM to reproduce the target functionality fully.
This bidirectional consistency check, combined with rigorous human verification (\cref{sec:human_eval}), ensures strict semantic alignment between the abstract requirements and the ground truth implementation.

\paragraph{\textbf{Reliability of Test-Based Evaluation.}}
The validity of our benchmark relies on whether passing the provided test suites truly indicates a correct and robust implementation.
First, to ensure high-quality evaluation, all test cases in FeatBench are curated directly from the original repositories.
These tests are authored by experienced human developers, ensuring a level of rigor and relevance that exceeds synthetic or auto-generated tests.
However, a potential concern remains about false positives, where incorrect code might pass a sparse suite by chance.
We mitigate this by enforcing a strict dual-validation strategy: candidate solutions must pass not only the newly added functional tests (F2P) to verify the feature but also a massive suite of existing regression tests (P2P).
This extensive coverage significantly raises the bar for success, ensuring that only solutions that are both functionally correct and free of regressions are considered resolved.

\paragraph{\textbf{Generalizability and Scope.}}
Currently, FeatBench focuses primarily on Python repositories.
While Python is a dominant language in modern software engineering, particularly in AI and data science, our findings may not fully extrapolate to statically typed languages like Java or C++.
However, the proposed benchmark construction methodology is language-agnostic and can be extended to other domains.
Furthermore, regarding data contamination, which is a critical threat in the LLM era, we mitigate this by designing FeatBench as an evolving benchmark.
By continuously sourcing tasks from the latest repository updates, we ensure that the evaluation data remains unseen during the training of the evaluated models.